\section{The interface with a renormalized quantum induction}
\label{gex-next-section}

The low-degree classification is now tied to the actual local
symmetries and a normalized classical reference. The exact
coupling rebase is
\[
 g_{\mathrm{next}}^{-2}=g^{-2}-\beta,\qquad
 E_{\mathrm{ren}}=E-E(0)-\beta S_{MR}.
\]
This is an identity of the density after extraction. A sign
or trajectory estimate for the generated \(\beta\) is still
needed; the identity does not supply one.

There is a concrete route to the spatial part of a contraction.
Fix the elementary macrostep \(R\) and its sufficient auxiliary
grouping \(J\). The spatial remainder estimate has constants
depending on these fixed data. Select a finite \(q\) afterward,
so that its endpoint ratio \(L=R^q\) makes
\(C(J)L^{-\alpha}\) small. The classical semigroup identifies
the scalar inherited background with the direct endpoint
background. Its new connected localization can then be
estimated directly. The high-degree piece has the separate
inverse-\(L\) estimate of Notebook 28, with its analytic radius
chosen after this finite \(L\).

The analytic compactness and quotient constructions make
this classical change of representation precise. They supply
joint limiting interactions and maps, a complete quotient,
and a bounded full-reference projection. The final classical
estimate is a contraction once the finite endpoint ratio
is sufficiently large and its common radius is chosen.

To turn this into a quantum induction, four interfaces must
be established. First, the limiting interaction space must
be matched to the regulators used by the actual quantum update,
including winding terms and its vacuum normalization.
Second, prove that the quantum update descends continuously
to the quotient; classical descent alone does not identify
the sharp finite-regulator integrals for equivalent inputs.
Third, the \(q\) successive
quantum steps must retain all intermediate generated terms,
cutoffs and normalizations; the classical semigroup does not
replace them by a single fluctuation integral. Fourth, the
resulting recurrence must control both the renormalized
activity and the actual coupling trajectory uniformly in
the number of blocks of steps.

The present results supply the nonlinear marginal identity,
the classical spatial gain and its representation space
under stated hypotheses. The larger existence proof
still needs the mixed and large-field histories, sourced
observable estimates, continuum nontriviality, Euclidean
symmetry and reconstruction already recorded in Notebook 01.
The mass gap remains a subsequent spectral assertion.

\begin{thebibliography}{9}
\bibitem{priorgex}
\emph{Toward a four-dimensional Yang--Mills construction},
research Notebooks 07, 15--16 and 25--28 in this workspace.
The flat classical normalization, ordinary minimum theorem,
strict classical support families and common-domain quantum
macrostep are the principal prior inputs.
\bibitem{balaban109gex}
T.~Ba\l aban,
\emph{Renormalization group approach to lattice gauge field
theories. I. Generation of effective actions in a small field
approximation and a coupling constant renormalization in four
dimensions},
Commun. Math. Phys. \textbf{109} (1987), 249--301.
\href{https://doi.org/10.1007/BF01215223}{Publisher record};
\href{https://scispace.com/pdf/renormalization-group-approach-to-lattice-gauge-field-3ygai0vxt1.pdf}{Original article copy}.
The averaging symmetries and total Ward/subtraction calculations
are checked against this primary source. Its stronger
coefficientwise orbit hypotheses are not assumed here.
\bibitem{balaban116gex}
T.~Ba\l aban,
\emph{Renormalization group approach to lattice gauge field
theories. II. Cluster expansions},
Commun. Math. Phys. \textbf{116} (1988), 1--22.
\href{https://doi.org/10.1007/BF01239022}{Publisher record}.
\bibitem{sourcegex}
E.~Sussman and collaborators,
\href{https://github.com/ethanwsussman/Yang--Mills-LLM-experiment}
{Yang--Mills-LLM-experiment},
deposited raw accounts at commit
\texttt{e392ef45013a581664bab76b90867170b30fe10e}.
\end{thebibliography}
\end{document}
